 % !TEX root = ../mainfile.tex
% !TEX encoding = UTF-8 Unicode
% !TEX spellcheck = en_US

\chapter{Literature Review}
\label{ch:literature}

This chapter situates the thesis within five strands of research---currency
factor models, safe-haven currencies and gold, geopolitical risk, tariff shocks,
and commodity currencies---and closes by distilling them into testable
hypotheses. The unifying question is when, and through which channel, the U.S.\
dollar functions as a safe haven, and whether that function survives when the
United States is itself the source of the shock.


\section{Currency Returns and Factor Models}
\label{sec:lit_currency_factors}

The modern cross-section of currency returns is organised around a small number
of risk factors. \textcite{lustig2011} identify a ``slope'' factor (the return on
a portfolio long high-interest-rate and short low-interest-rate currencies) that
explains most of the cross-sectional variation in carry returns and is tied to
innovations in global equity-market volatility; high-interest-rate currencies
load more heavily on this common factor, so investors in the carry trade are
compensated for bearing global risk. This logic motivates sorting currencies into
portfolios by their exposure to common shocks, and the dollar and carry factors
derived from it underpin the portfolio construction used in
Chapter~\ref{ch:methodology}. Related work on dominant-currency pricing links the
pricing of trade and the currency composition of invoicing to these risk premia
\parencite{dalgic2024,dalgic2025}.

The dollar occupies a special position in this system. \textcite{gourinchasrey2005}
document the United States' ``exorbitant privilege''---a persistent excess return
on its external assets over its liabilities---and \textcite{gourinchas2017} show
that this privilege is mirrored by an ``exorbitant duty'': the hegemon transfers
wealth to the rest of the world precisely during global crises, with transfers of
roughly $19\%$ of U.S.\ GDP during the 2007--2009 financial crisis. \textcite{gourinchas2019}
embed this in a broader account of the international monetary system and a renewed
Triffin dilemma, while \textcite{obstfeld2023} describe a ``global dollar cycle''
in which dollar appreciation tightens financial conditions worldwide. Together these
contributions frame the dollar's safe-haven status as a structural feature of the
international monetary system---and therefore something whose erosion would be
economically consequential.


\section{Safe-Haven Currencies and Gold}
\label{sec:lit_safe_haven}

Following \textcite{baur2010}, an asset is a \emph{hedge} if it is uncorrelated
with risky assets on average and a \emph{safe haven} if it is uncorrelated with
(or negatively correlated with) them specifically during market stress; they find
gold to be a hedge against equities on average and a safe haven in extreme
declines, though the safe-haven property is short-lived. This conditional, crisis-
contingent definition is the one adopted throughout the thesis. Applied to
currencies, \textcite{ranaldo2010} show that the Swiss franc and Japanese yen
appreciate against the dollar when U.S.\ equities fall and volatility rises, with
the effect strengthening non-linearly in the volatility factor and during crises.
\textcite{debock2015} document the same recurrent pattern across risk-off
episodes---the yen, franc and dollar all appreciate against other currencies---and
relate it to fundamentals such as the nominal interest rate and the net foreign
asset position. \textcite{habib2012} push beyond the carry trade to ask what
\emph{makes} a currency a safe haven, finding that the net foreign asset position,
low external vulnerability, and economic size are robustly associated with
safe-haven status, whereas the interest-rate spread against the United States is
not. Methodologically, the time-varying correlation approach used to test these
properties---and applied to gold and other assets by \textcite{kumar2025}---rests
on the dynamic conditional correlation model discussed in Section~\ref{sec:lit_gpr}
and Chapter~\ref{ch:methodology}.


\section{Geopolitical Risk and Currency Markets}
\label{sec:lit_gpr}

\textcite{caldara2022} construct the news-based Geopolitical Risk (GPR) index used
in this thesis, showing that it spikes around wars and major adverse events and
that higher geopolitical risk foreshadows lower investment and employment and
larger downside risks; importantly, they distinguish the \emph{threat} of adverse
events from their \emph{realisation}, a distinction directly relevant to staged
crises such as the Hormuz episode. Building on such measures, \textcite{liu2024}
study geopolitical-risk-sorted currency portfolios and report a positive risk
premium for exposure to geopolitical shocks, and the \textcite{imf2025gfsr} finds
that geopolitical-risk shocks predict currency depreciation and equity declines,
particularly in more vulnerable economies. This literature motivates treating
geopolitical events as priced risk rather than noise, and supplies the GPR series
used as an external control in the empirical chapters.


\section{Tariff Shocks and Exchange Rates}
\label{sec:lit_tariffs}

The 2025 tariff episode has generated a fast-growing literature that is central to
this thesis. \textcite{kamin2025} provides econometric evidence that, following
the Liberation Day announcements of April~2025, the dollar switched from behaving
as a safe haven---appreciating in periods of market volatility---to behaving as a
``risk-on'' currency that moves inversely with volatility. \textcite{jiang2026}
document the same break in historical correlations (the dollar depreciating while
U.S.\ rates rose relative to the euro, the VIX increased, and the dollar
convenience yield fell) and interpret it through shifts in global demand for
dollar safe assets, calibrating a steady-state real dollar depreciation of around
$7.6\%$ should the United States lose reserve-currency demand. \textcite{hassan2025}
provide the theoretical counterpart: in a general-equilibrium model where currency
safety is endogenous, a trade war that isolates U.S.\ goods markets erodes the
dollar's safety premium and raises U.S.\ interest rates. On the empirical
exchange-rate response, \textcite{corsetti2025} construct a tariff-shock database
and show that exchange rates react to U.S.\ tariff shocks systematically once the
size and relevance of the measures (and expected retaliation) are accounted for,
while \textcite{kalemliozcan2026} emphasise the role of tariff \emph{uncertainty}
in the dollar's behaviour. The theoretical tariff--exchange-rate offset is
formalised by \textcite{itskhoki2025}. Policy institutions corroborate the market
dislocation: the \textcite{ecb2025fsr} and the \textcite{ecb2026bulletin} document
trade fragmentation, the pass-through of higher U.S.\ tariffs, and the effects of
trade-policy uncertainty on activity. This thesis contributes a direct event-study
comparison of this trade-policy shock against contemporaneous military/oil-supply
shocks.


\section{Commodity Currencies and Oil}
\label{sec:lit_commodities}

\textcite{chen2003} establish that the U.S.-dollar price of commodity exports
exerts a strong and stable influence on the real exchange rates of commodity
exporters such as Australia, Canada and New Zealand, providing the basis for the
oil-exporter versus oil-importer portfolio split used in
Chapter~\ref{ch:methodology}. The nature of the oil shock matters, however:
\textcite{kilian2009} shows that oil-price movements driven by supply
disruptions, aggregate demand, and precautionary (oil-specific) demand have
distinct macroeconomic consequences, so a price increase caused by a supply scare
is not equivalent to one caused by strong demand---nor to a tariff-induced demand
collapse. This typology maps directly onto the thesis's contrast between the
demand-side tariff shock (falling oil) and the supply-side Hormuz shock (surging
oil). For the 2026 episode specifically, \textcite{fattouh2026} provide a detailed
anatomy of the Strait of Hormuz oil shock and the associated disruption to Gulf
output, while \textcite{kilian2026} quantify the inflationary consequences of the
2026 Iran war through oil prices, and the \textcite{imf2026menap} documents the
regional and global spillovers, including Brent rising above \$100 per barrel and
the role of the ceasefire in de-escalation.


\section{Hypotheses}
\label{sec:hypotheses}

Taken together, these strands point to a single organising prediction: the dollar's
safe-haven response is conditional on the \emph{type} of shock rather than on whether the
United States is its originator. Both the external benchmark (the 2022 Russian invasion of
Ukraine) and the U.S.-originated Hormuz crisis are military/oil-supply shocks, around which
the safe-haven literature (Section~\ref{sec:lit_safe_haven}) implies the dollar should
appreciate; the Liberation Day tariff announcement, by contrast, is a U.S.\ trade-policy
shock that the dollar-erosion literature (Section~\ref{sec:lit_tariffs}) suggests should
weaken the very currency it originates from. Gold offers a natural counterpoint: as an asset
external to any sovereign it might be expected to bid in \emph{both} kinds of episode, so
that comparing it against the dollar isolates what is specific to the dollar's status
\parencite{baur2010}---a prediction this thesis puts directly to the data. The cross-section
of currencies should in turn sort on the oil channel, exporter currencies outperforming
importers when the Hormuz supply shock lifts crude but losing ground when the tariff shock
depresses oil through weaker demand \parencite{chen2003,kilian2009}.

\noindent The remainder of the thesis puts these predictions to the data: the dollar
and the oil cross-section in the event study of Chapter~\ref{ch:event_study}, with gold
examined alongside them, while a regression test in Chapter~\ref{ch:regression} establishes
whether the dollar's safe-haven relationship with risk held or broke in each episode.
